\input{../tex-inputs/latex-leseansicht-vorspann}

\section[Arthur und Olga Schnitzler an Gustav Schwarzkopf, 16. 8. 1906]{L04381 Arthur und Olga Schnitzler an Gustav Schwarzkopf, 16. 8. 1906}
\nopagebreak\mylabel{L04381v}
\rehead{ }\normalsize\beginnumbering\briefempfaengerindex{Schwarzkopf, Gustav@\textsc{Schwarzkopf, Gustav}!zzzSchnitzler, Olga@\emph{von Olga Schnitzler}!1906-08-161@{16. 8. 1906}|(be}\briefempfaengerindex{Schwarzkopf, Gustav@\textsc{Schwarzkopf, Gustav}!zzzSchnitzler, Arthur@\emph{von Arthur Schnitzler}!1906-08-161@{16. 8. 1906}|(be}
\toendnotes[C]{\smallbreak\pagebreak[2]}
\correspDesc{Versand  durch Arthur Schnitzler, Olga Schnitzler am 16. 8. 1906 in Ilmenau
\newline{}Übermittlung  am 16. 8. 1906 in Ilmenau
\newline{}Erhalt  durch Gustav Schwarzkopf im Zeitraum [17. 8. 1906 – 21. 8. 1906?] in Wien}\toendnotes[C]{\smallbreak}
\Standort{DLA, A:Schnitzler, HS.1985.1.1897.}
\physDesc{Bildpostkarte, 173 Zeichen
\newline{}Handschrift Arthur Schnitzler: Bleistift, deutsche Kurrent
\newline{}Handschrift Olga Schnitzler: Bleistift
\newline{}Versand: Stempel: »\nobreak{}\oindex{Ilmenau@\textbf{Ilmenau}|pwk}Ilmenau, 16. 8. 1906, 11–12N\nobreak{}«.  }\toendnotes[C]{\smallbreak}\pstart{}{\pb}Hrn Guſtav Schwarzkopf\pend{}\pstart{}Wien I\oindex{I., Innere Stadt@\textbf{I., Innere Stadt}|pw}\pend{}\pstart{}Tiefer Graben 23\oindex{Wien@\textbf{Wien}!I., Innere Stadt@\textbf{I., Innere Stadt}!Tiefer Graben 23@\textbf{Tiefer Graben 23}|pw}.\pend{}{\bigskip}
\pstart
           \centering{}{\pb}\textcolor{gray}{\textbf{Göthehäuschen\oindex{Goethehäuschen@\textbf{Goethehäuschen}|pw} bei Bad Ilmenau\oindex{Ilmenau@\textbf{Ilmenau}|pw}.}}\pend
           \vspace{1em}
\pstart
           \noindent{}{\pb}Herzliche Grüße! Wo ſind Sie? Wir hoffen \label{K_L04381-1v}\edtext{in acht Tagen (nach Eiſenach\oindex{Eisenach@\textbf{Eisenach}|pw} u Nürnberg\oindex{Nürnberg@\textbf{Nürnberg}|pw}) zu
                  Haus}{\lemma{\textnormal{\emph{in … Haus}}}\Cendnote{\textnormal{Sie kamen am 21. 8. 1906 und damit
                  sogar früher als gedacht nach Wien\oindex{Wien@\textbf{Wien}|pwk}.}}}\label{K_L04381-1}
               zu ſein.\pend
           \pstart \spacefill\mbox{Arthur}\pend{}
\pstart
           \noindent{}Schönen Gruß an Dr. Max\pwindex{Schwarzkopf, Max 12.\,6.\,1857 Wien – 14.\,4.\,1928 ebd.@\textsc{Schwarzkopf, Max} (12.\,6.\,1857 Wien – 14.\,4.\,1928 ebd.)|pw}.\pend
           \selectlanguage{ngerman}\vspace{1em}\pstart \spacefill\mbox{{[}hs. O. Schnitzler:{]} O. S.}\pend{}\selectlanguage{ngerman}\endnumbering\briefempfaengerindex{Schwarzkopf, Gustav@\textsc{Schwarzkopf, Gustav}!zzzSchnitzler, Olga@\emph{von Olga Schnitzler}!1906-08-161@{16. 8. 1906}|)be}\briefempfaengerindex{Schwarzkopf, Gustav@\textsc{Schwarzkopf, Gustav}!zzzSchnitzler, Arthur@\emph{von Arthur Schnitzler}!1906-08-161@{16. 8. 1906}|)be}\mylabel{L04381h}
\begin{anhang}
\end{anhang}\newcommand{\dateiname}{L04381}\newcommand{\titel}{Arthur und Olga Schnitzler an Gustav Schwarzkopf, 16. 8. 1906}\newcommand{\editorInnen}{Herausgegeben von Jahnke, SelmaMüller, Martin Anton}\input{../tex-inputs/latex-leseansicht-abspann}
